\documentclass[11pt]{article}
\usepackage[margin=1in]{geometry}
\usepackage{amsmath,amssymb,booktabs}
\usepackage[hidelinks]{hyperref}
\title{Exact certificates for sampled 17-city subtour vertices}
\author{Claude Make No Mistakes}
\date{25 November 2026}
\begin{document}
\maketitle

\begin{abstract}
We sampled 24 labeled non-half-integral vertices of the 17-city subtour
polytope and certified the exact worst metric ratio for each. The largest
ratio in the sample was \(15/13\). A 26-edge support gave \(10/9\).
The search produced no counterexample to the metric TSP \(4/3\) conjecture.
\end{abstract}

\section{Certificate}
Let \(V\) have \(n\geq3\) cities and \(E=\binom{V}{2}\). The subtour polytope
consists of \(x\in[0,1]^E\) with \(x(\delta(v))=2\) at each vertex and
\(x(\delta(S))\geq2\) for every nonempty proper \(S\subset V\).
For a complete metric \(c\) let \(T(c)\) denote the minimum Hamiltonian
tour cost. For a fixed feasible \(x\), define its worst metric ratio as
\[
    g(x)=\sup_{c\text{ a metric on }V} \frac{T(c)}{c\cdot x}.
\]
The ordinary integrality gap is \(\sup_x g(x)\). For each metric, the
denominator is minimized by an extreme point, so surveying extreme
points targets the conjecture directly \cite{chp}.

For a lower bound, assign nonnegative integer lengths \(w\) to
\(G=(V,\operatorname{supp}x)\), with \(w\cdot x>0\). Let \(d\) be the
shortest-path distances in \(G\). Feasibility of \(x\) ensures connectivity.
Every supported edge satisfies \(d_e\leq w_e\). If \(d\) has zero distances
between distinct cities, add \(\varepsilon>0\) to each off-diagonal
distance to obtain a metric \(d_\varepsilon\). Since \(x(E)=n\) and each
tour has \(n\) edges,
\[
    T(d_\varepsilon)=T(d)+n\varepsilon,\qquad
    d_\varepsilon\cdot x\leq w\cdot x+n\varepsilon.
\]
Taking \(\varepsilon\) down to zero gives \(T(d)/(w\cdot x)\leq g(x)\).
A strict lower bound above \(4/3\) would give a metric counterexample.

For an upper bound, take connected spanning Eulerian multigraphs on
\(G\) with edge multiplicities \(a^j\) and rational probabilities
\(p_j>0\) summing to one. By shortcutting an Euler tour in any metric
\(c\), \(T(c)\leq c\cdot a^j\). If their expected edge multiplicities
are bounded by \(Ux\) coordinatewise, then
\[
    T(c)\leq c\cdot\sum_j p_j a^j\leq U(c\cdot x).
\]
Matching lower and upper bounds determine \(g(x)\) exactly.

\newpage
\section{Generation and verification}
We start from two triangles with corresponding vertices joined by
internally disjoint paths of lengths \(4,5,5\). The triangle edges
have weight \(1/2\) in \(x\), and the path edges have weight \(1\).
We add random chords and optimize an objective on the resulting
restricted subtour polytope by cut separation. We retain reconstructed
points with at least 24 support edges and at least one coordinate
outside \(\{0,1/2,1\}\). The three runs change the available chords or
bias the objective toward the positive control's metric.

For each retained point we search for lengths \(w\) minimizing
\(w\cdot x\) subject to each discovered spanning Eulerian walk having
length at least 1. A shortest tour in the graph metric closure expands
to a separating walk. Approximate LP duals give candidate mixtures.
The certificate includes integer lengths and rational weights. Its
validity is independent of floating-point solver tolerances.

The verifier uses exact flows from city 0 to every other city to
certify cuts, and exact Gaussian elimination on degree equations
and tight constraints to certify vertex status. All absent edges
are fixed at zero. It also checks connectivity and even degrees
of all multigraphs and recomputes \(T(d)\) with a separate integer
Held--Karp implementation. Results are bundled as JSON Lines.

\begin{center}
\begin{tabular}{lrrr}
\toprule
Search & Attempts & Labeled vertices & Largest \(g(x)\) \\
\midrule
Unrestricted chords & 1600 & 20 & \(23/20\) \\
Restricted chord endpoints & 1800 & 2 & \(15/13\) \\
Control-metric bias & 2200 & 2 & \(129/113\) \\
\bottomrule
\end{tabular}
\end{center}

The positive control yielded \(53/43\), as reported for the
17-city half-integral maximum in \cite{chp}. The sample includes
15 vertices with 24 support edges, eight with 25, and one with 26.
The search draws from structured supports and provides conclusions
only for the recorded vertices.

\section{Relation to existing results}
Cook, Hougardy, and Petrich verify the conjecture through 16 cities
in general and for half-integral vertices through 18 cities \cite{chp}.
Villa, Vercesi, Barta, and Mastrolilli prove the \(n+6\) support bound
\cite{vvbm}. Public artifacts report extensions to \(n+7\) and \(n+8\)
\cite{regev,naser}. The 26-edge case lies beyond those support limits.

\begin{thebibliography}{9}
\bibitem{chp} W. Cook, S. Hougardy, and M. Petrich.
\emph{Extending exact integrality gap computations for the metric TSP}.
\href{https://arxiv.org/html/2603.12995v5}{arXiv:2603.12995v5}, 2026.
\bibitem{vvbm} T. Villa, E. Vercesi, J. Barta, and M. Mastrolilli.
\emph{The integrality gap of the traveling salesman problem is \(4/3\)
if the LP solution has at most \(n+6\) non-zero components}.
\href{https://arxiv.org/abs/2507.07003v3}{arXiv:2507.07003v3}, 2026.
\bibitem{regev}
\emph{An exact certificate proof for the \(n+7\) metric TSP support case}.
\href{https://github.com/GiladR1979/metric-tsp-nplus7-certificates}{Public proof artifact}, 2026.
\bibitem{naser} H. M. Naser.
\emph{The metric TSP \(4/3\) bound with support at most \(n+8\):
Structural reduction and slice-first enumeration}.
\href{https://ssrn.com/abstract=7395957}{SSRN 7395957}, 2026.
\end{thebibliography}
\end{document}
